\documentclass{book}
\usepackage[utf8]{inputenc}
\usepackage[T1]{fontenc}
\usepackage{geometry}
\usepackage{graphicx}
\usepackage{float}
\usepackage{hyperref}
\usepackage[english,indonesian]{babel}

\title{Optimasi Pemrosesan Template LaTeX Menggunakan Pendekatan Berbasis AI}
\author{Academic Support Agent}
\date{\today}

\begin{document}
\maketitle

\begin{abstract}

\end{abstract}

\chapter{PENDAHULUAN}


\chapter{TINJAUAN PUSTAKA}


% Original content (commented out):
% % University of Strathclyde
% Research dissertation template
% for degrees from the Mechanical \& Aerospace Engineering
%------------------------------------------------
\documentclass[a4paper,oneside,12pt]{book}


% fonts
\usepackage{times} 
\usepackage{amsbsy, amsmath, amsfonts}
\usepackage{siunitx}
\usepackage{graphicx}

% text sub- and superscripts
\newcommand{\spr}[1]{\ensuremath{^\textrm{#1}}}
\newcommand{\sub}[1]{\ensuremath{_\textrm{#1}}}

% better tables
\usepackage{multirow}
\usepackage{booktabs}

% links
\usepackage{color}
\usepackage[]{hyperref}
\usepackage{enumitem} % better lists

% Use for bibliography style, Harvard or Vancouver, etc.
\usepackage{natbib}

% figures and tables
\usepackage{caption}
\captionsetup[figure]{labelfont=it,textfont=it}
\captionsetup[table]{labelfont=it,textfont=it}
\usepackage{subfig}

\usepackage[strict]{changepage}

% Redefine \eqref to include Eq. prefix
\makeatletter
\def\tagform@#1{\maketag@@@{(\ignorespaces#1\unskip\@@italiccorr)}}
\makeatother
\let\originaleqref=\eqref
\renewcommand{\eqref}{Eq.~\originaleqref}

% Define margins
\usepackage[left=2.5cm,right=2.5cm,top=2cm,bottom=4cm,includehead,includefoot,headheight=15pt]{geometry}

% Page headers
\usepackage{fancyhdr}
\fancyhf{}
\renewcommand{\headrulewidth}{0pt} % optional
%\fancyhead[L]{\nouppercase{\leftmark} \hfill Section \nouppercase{\rightmark}}
\fancyhead[L]{\nouppercase{\leftmark}}
\cfoot{\thepage}
\pagestyle{fancy}


% Line spacing
%\def\baselinestretch{1.5} 
\usepackage{setspace}
\setstretch{1.5}

% Defines strath-specific title page
\usepackage{strathtitle}

% Nomenclature
\usepackage{nomencl}
\makenomenclature

\usepackage{etoolbox}  % option to group symbols
\renewcommand\nomgroup[1]{%
  \item[\bfseries
  \ifstrequal{#1}{C}{Constants}{%
  \ifstrequal{#1}{S}{Variables}{}}%
]}

%-------------------------------------------
\begin{document}

% Complete with your own details
\author{First Surname}
\title{Title}
\course{NAME of THE COURSE/DEGREE }
\supervisor{Dr First Surname}

\maketitle

\frontmatter

\chapter*{Abstract}
Abstract should be on one page. It should summarise the project objectives, the work carried out, methods used, the main research findings arising from the work and conclusions reached. 

\chapter*{Acknowledgements}
This thesis is the result of the author’s original research. However if any external support, help or contributions were received, it must be acknowledged here. 

\tableofcontents

\addcontentsline{toc}{chapter}{List of Figures}
\listoffigures

\addcontentsline{toc}{chapter}{List of Tables}
\listoftables

\addcontentsline{toc}{chapter}{Nomenclature}
\mbox{}
\nomenclature[C]{$c$}{Speed of light in a vacuum, \SI{299792458}{m/s}}
\nomenclature[C]{$h$}{Plank's constant, \SI{6.62607004e-34}{m^2.kg/s}} %m\spr{2}kg/s
\nomenclature[S]{$g$}{Gravitational acceleration, m/s\spr{2}}
\nomenclature[S]{$\mathbf{q}$}{Quaternions}
\nomenclature[S]{$V$}{Volume, m\spr{3}}
\nomenclature[S]{$\rho$}{Density, kg/m\spr{3}}

\printnomenclature

%-------------------------------------------------
\mainmatter

\chapter{Introduction}

\section{Format of Tables, Figures and Equations}
Each item must have a numerical label and title (a caption). Figure captions must appear below the figure, see Fig.\ \ref{fig:wave}. Table captions must appear above the table, see Table \ref{tab:my_label}. Equation captions appear to the right of the equation (or below if space does not permit), see \eqref{eq:y}.

A decaying harmonic is given as,
\begin{equation}\label{eq:y}
    y(t)=a e^{-b\omega t} \sin(\omega t + \phi)
\end{equation}
where $a$ is the sine wave amplitude, frequency $\omega$, phase offset $\phi$, and the exponential coefficient $b>0$. An example is shown in Figure \ref{fig:wave}. 

\begin{figure}[thb]
	\centering
	\includegraphics[width=0.9\textwidth]{figure_wave.eps}
	\caption{Exponentially decaying sine wave from \eqref{eq:y}}\label{fig:wave}
\end{figure}

This is generally an infinite dimensional optimisation problem of the form
\begin{equation}\label{eq:opt}
\min\limits_{u,x} \;\phi\left(x(t_f)\right) \;+\int_{t_0}^{t_f} f_0\left(x(t),u(t)\right)\,dt
\end{equation}
subject to
\begin{align}
\dot{x}(t)&=f\left(x(t),u(t)\right)\\
c\left(x(t),u(t)\right) &\leq 0,\;\; t\in[t_0,t_f]\\
\omega\left(x(t_0),x(t_f)\right) &= 0 
\end{align}
where $f(x(t), u(t))$ describes the dynamic of the system, the function $\omega(x(t_0),x(t_f))$ the equality constraints on initial and final states and $c(x(t),u(t))$ the path constraints.
The solution of an optimal control problem is the pair $[x^*(t),u^*(t)]$ which minimise the objective function defined by $\phi(x(t_f))$ and $f_0(x(t),u(t))$.

Tables, Figures and Equations must have separate numbering. Each item should appear in sequential order and be mentioned explicitly in the body of text, and be mentioned prior to appearing in the report. If you are using a figure or table from a third party (book, paper, website etc.) then it must be referenced in the caption, as in Tables \ref{tab:my_label}--\ref{tab:my_label2}. If data has been extracted from several sources and compiled into a table, this can be referenced in the table itself.

Each item should be centralised in the body of text. No text wrapping should be employed around figures, tables, or equations. All text in the figure should be legible. There is no limitation of figure size, it is better to show a figure larger and legible, than smaller and indecipherable.

\begin{table}[htb]
    \centering
    \begin{tabular}{l l l p{5cm}}
    \toprule
    Day & Min Temp & Max Temp & Summary \\ \midrule
    Monday & \SI{11}{\celsius} & \SI{22}{\celsius} & A clear day with lots of sunshine.  
    However, the strong breeze will bring down the temperatures. \\ \midrule
    Tuesday & \SI{9}{\celsius} & \SI{19}{\celsius}  & Cloudy with rain, across many northern regions. Clear spells 
    across most of Scotland and Northern Ireland, 
    but rain reaching the far northwest. \\ \midrule
    Wednesday & \SI{10}{\celsius}  & \SI{21}{\celsius}  & Rain will still linger for the morning. 
    Conditions will improve by early afternoon and continue 
    throughout the evening. \\
    \bottomrule
    \end{tabular}
    \caption{Table caption \citep{website:table}}
    \label{tab:my_label}
\end{table}

\begin{table}[htb]
    \centering
\begin{tabular}{llr}  
\toprule
\multicolumn{2}{c}{Item} \\
\cmidrule(r){1-2}
Animal    & Description & Price (\$) \\
\midrule
Gnat      & per gram    & 13.65      \\
          &    each     & 0.01       \\
Gnu       & stuffed     & 92.50      \\
Emu       & stuffed     & 33.33      \\
Armadillo & frozen      & 8.99       \\
\bottomrule
\end{tabular}
    \caption{Caption for table taken from \cite{website:table}}
    \label{tab:my_label2}
\end{table}

\section{Nomenclature}
All variables and symbols used should be given in the text, and summarised in List of Symbols, or Nomenclature at the start. 
The nomenclature should include all the variables, their descriptions and if applicable, the values (in the case of constants) and units. It does not usually include acronyms. 

The symbols can be grouped, or not, to help with clarity depending on the number of variables.


\section{Referencing}
Referencing should preferably follow the \href{https://library.aru.ac.uk/referencing/harvard.htm}{Harvard system}, although a numeric style such as \href{https://en.wikipedia.org/wiki/Vancouver_system}{Vancouver system} is also acceptable. It is important that attention is paid to proper referencing of others' work (including figures), both to avoid plagiarism and to allow others to find the source document \citep{blackholes}, e.g., web links should contain the URL and the date that the web link was sourced. If necessary, Turnitin will be used to check for plagiarism.

A bibliography can be used to reference different reference sources including books \citep{book:goossens93}, and articles \citep{art:AbedonHymanThomas2003,art:greenwade93}. \citet{proc:Abedon1994} wrote an interesting paper in the field of molecular biology, which has a slightly different citation format.

%--------------------------------------------
\chapter{Literature Review}
\section{Section 1}

\section{Section 2}
\subsection{Sub-section}

\section{Section 3}

%--------------------------------------------
\chapter{Method}


%--------------------------------------------
\chapter{Results}


%--------------------------------------------
\chapter{Discussion}



%--------------------------------------------
\chapter{Conclusion}
This section should summarise the information presented in the report. No new information regarding the work should be included here. It will often include a further or future work section, discussing what would be useful going forward.


%----------------------------------------------
% Inserts references/bibliography, using myrefs.bib as a database but only includes those entries referenced in the text above
\addcontentsline{toc}{chapter}{References}
\renewcommand{\bibname}{References}
\bibliographystyle{agsm} % Harvard style
\bibliography{myrefs} 

%-----------------------------
\appendix
\chapter{Additional material}
Appendices give the chance to include supplementary information which is useful but not pertinent to the understanding of the report. Appendices should be referred to explicitly in the main text, or should not be included.	The appendices shall not be used as a means to include additional word count.



\end{document}

\end{document}